\section{Từ lời giải thích cục bộ đến một kết luận toàn cục}
\label{sec:ch04-tu-cuc-bo-den-toan-cuc}

Vì TreeSHAP chỉ làm phép tính cục bộ nhanh hơn, một kết luận toàn cục vẫn phải
được xây từ nhiều đầu vào. Mỗi vector $\phi_i(x)$ giải thích một dự đoán tại $x$. Trong thực hành, người
ta thường lấy trung bình độ lớn hoặc bình phương của các giá trị này trên một
tập dữ liệu để xếp hạng đặc trưng. Phép tổng hợp ấy chuyển một lời giải thích
cục bộ thành một phát biểu toàn cục, nên nó cần một câu hỏi mới: độ bất định
của con số tổng hợp là bao nhiêu?

Whitehouse, Sawarni và Syrgkanis xét các lũy thừa của đường cong SHAP, trong
đó có trung bình trị tuyệt đối và trung bình bình phương. Họ chỉ ra rằng việc
cắm một mô hình đã học vào rồi lấy trung bình đơn giản thường không đủ để có
suy luận tiệm cận chuẩn; bài báo xây dựng các ước lượng đã hiệu chỉnh để lập
khoảng tin cậy cho các đại lượng đó~\autocite{p19shapinference}.

Kết quả này không biến một giá trị SHAP cục bộ thành bằng chứng về cơ chế của
mô hình. Nó làm một việc khác: nếu ta đã chọn hàm giá trị và dữ liệu đánh giá,
thì phát biểu toàn cục dựa trên các giá trị SHAP phải mang theo bất định lấy
mẫu và bất định do ước lượng. Đây là nền tảng thống kê mà một biểu đồ xếp hạng
đặc trưng thường chưa thể hiện.

\index{suy luận thống kê}%
\index{lời giải thích cục bộ}%
